\section{Titus Flavius Clemens and the Golden Chain}

What evidence could we adduce for the claim that Clement of Alexandria possessed secret wisdom?

First of all, we could establish his difference from (say) Tertullian or even Jerome (who had the dream about getting rid of his pagan library):

\begin{quotex}
... by philosophy I mean not the Stoic, nor the Platonic, nor the Epicurean, nor that of Aristotle; but whatever any of these sects had said that was fit and just, that taught righteousness with a divine and religious knowledge, this I call eclectic philosophy.
\end{quotex}


For Clement, Christianity was true, not because it overthrew paganism, but because it completed it. Furthermore, truth was truth, not ``Christian'' truth. Yet he goes somewhat further, differing from (say) Augustine, who first believed in order to understand, writing instead that

\begin{quotex}
As many men drawing down the ship, cannot be called many causes, but one cause consisting of many; --- for each individual by himself is not the cause of the ship being drawn, but along with the rest; --- so also philosophy, being the search for truth, contributes to the comprehension of truth; \textbf{not as being the cause of comprehension, but a cause along with other things, and co-operator; perhaps also a joint cause}. ... But the Hellenic truth is distinct from that held by us (although it has got the same name), both in respect of extent of knowledge, certainly of demonstration, divine power, and the like. For we are taught of God, being instructed in the truly ``sacred letters'' by the Son of God. \textbf{Whence those, to whom we refer, influence souls not in the way we do, but by different teaching.} And if, for the sake of those who are fond of fault-finding, we must draw a distinction, by saying that philosophy is a concurrent and cooperating cause of true apprehension, being the search for truth, then we shall avow it to be a preparatory training for the enlightened man (\emph{to gnostikon}); not assigning as the cause that which is but the joint-cause; nor as the upholding cause, what is merely co-operative; \emph{nor giving to philosophy the place of a sine qua non.} Since almost all of us, without training in arts and sciences, and the Hellenic philosophy, and some even without learning at all, through the influence of a philosophy divine and barbarous, and by power, have through faith received the word concerning God, \textbf{trained by self-operating wisdom}. But that which acts in conjunction with something else, being of itself incapable of operating by itself, we describe as co-operating and concausing, and say that it becomes a cause only in virtue of its being a joint-cause, and receives the name of cause only in respect of its concurring with something else, but that it cannot by itself produce the right effect.
\end{quotex}


The reserve about ``sine qua non'' is perhaps worth further debate. However, one understands, in order to believe. Yet one might go (in fact) further than this and construe that for Clement, Christianity could only remain true to the extent that it was also willing to be completed in its turn, or re-completed. That is, what took absolute precedence over exoteric doctrine would have been the divine Logos which from the beginning undertook the creation of the world as a school for the souls of men, a school which was living and active, creative and artful, not capable of being uttered in fullness in a creed. This plan would have been understood as articulated (and articulate) through initiation into what was pellucid, beautiful, and virtuous, the vestibules of the holy. ``He who finds them, finds more.'' The mastery of such would have been gnosis, the goal of pistis (only in this sense inseparable, but absolutely separate otherwise, since Clement would never have separated pistis completely from gnosis, but defended the ``higher knowledge''):

\begin{quotex}
It is then, as appears, the greatest of all lessons to know one's self. For if one know himself, he will know God; and knowing God, he will be made like God... But that man with whom the Word dwells does not alter himself, does not get himself up: he has the form which is of the Word; he is made like to God... and that man becomes God, since God so wills. Heraclitus, then, rightly said, ``Men are gods, and gods are men.'' (ANF 2.271). \flright{\emph{The Instructor}, 3.1}

On this wise it is possible for the true Gnostic already to have become God. ``I said, Ye are gods, and sons of the highest.'' (ANF 2.437). \flright{\emph{Strom.}, 4.23}

By thus receiving the Lord's power, the soul studies to be God; ... To the likeness of God, then, he that is introduced into adoption and the friendship of God, to the just inheritance of the lords and gods is brought; if he be perfected, according to the Gospel, as the Lord himself taught. (ANF 2.506) \flright{\emph{Strom.}, 6.14}

... they are called by the appellation of gods, being destined to sit on thrones with the other gods that have been first put in their places by the Saviour. (ANF 2.539). \flright{\emph{Strom.}, 7.10}

What, then, shall we say of the true Gnostic himself? ``Know ye not'', says the apostle, ``that ye are the temple of God?'' The true Gnostic is consequently divine, and already holy, God-bearing and God-borne. (ANF 2.547). \flright{\emph{Strom.}, 7.13}

But he who has returned from this deception, on hearing the Scriptures, and turned his life to the truth, is, as it were, from being a man made a god. (ANF 2.551) \flright{\emph{Strom.}, 7.16}
\end{quotex}


Yet there is even more than demonstrating the fitness of Clement's teaching as a vessel for esotericism. Although the Mar Saba letter of Clement mentioning the ``secret Gospel of Mark'' is debated among scholars, we should not be entirely convinced by scholars who argue\footnote{\url{http://ixoyc.net/data/Fathers/622.pdf}} that Clement's lack of further quotation implies disapproval.

Clement himself held a position that truth was incommunicable at high levels, except by initiation, or proper coding. This is consistent with his Pythagoreanism\footnote{\url{http://www.bu.edu/wcp/Papers/Anci/AnciAfon.htm}}.

This is drawn out at length by S. Muller\footnote{\url{http://stephanhuller.blogspot.com/2011/01/how-to-read-and-understand-clement-of_23.html}}:

\begin{quotationx}
As we noted in our last post Clement used the conflict with regards to Paul accepting circumcision in Acts and the Epistle to the Galatians to illustrate `a hidden harmony' which is necessary for `right conduct' and `right belief.' Now Clement has moved beyond this and opened the door to things written in two different forms in two different `forms' of the gospel. What Jesus spoke in a low voice (whispered) are the parables which appear in the revealed gospels of Christianity which circulate generally. These parables represent a `lower' or subordinate form of the mystic wisdom ultimately revealed by Jesus and preserved in some other gospel form. Yet the harmony between these two gospel forms is indeed what Clement understands as the `ecclesiastic canon' -- a Pythagorean term which clearly represents again the preserving of an original tune in a `lower key' (perhaps distinguished from the tonic by a ``\emph{diatessaron}''). 

When we go back to Clement's original discussion it is important to see the underlying Pythagorean coloring to his use of the term \emph{kanona}. For we read again: ``He spake all things in parables, and without a parable spake He nothing unto them;'' and if ``all things were made by Him, and without Him was not anything made that was made,'' consequently also prophecy and the law were by Him, and were spoken by Him in parables. ``But all things are right,'' says the Scripture, ``before those who understand,'' that is, those who receive and observe, according to the exposition of the Scriptures explained by Him according to the ecclesiastic canon; and the ecclesiastical rule is the concord and harmony of the law and the prophets in the covenant delivered at the coming of the Lord.

Thus in no uncertain terms then, Clement understands the term `canon' in a Pythagorean sense. In this particular case Clement uses Jesus's condemnation of the Pharisees who `lie' about the truth of the Law and the prophets and its prediction of a messiah like Christ as an example of unworthy behavior in the temporary Alexandrian community. Whoever Clement is attacking (the Carpocratians? the Marcionites? both?) have divorced the secret gospel from its original `harmony' with the public gospels so as to ruin the true ecclesiastic `canon' understood in the Pythagorean sense.
\end{quotationx}


Salvatore Lilla\footnote{\url{http://journals.cambridge.org/action/displayAbstract;jsessionid=C80CC407F4BFC4A35AD9963BC675DA26.tomcat1?fromPage=online\&aid=2407972}} is agreed to have done good work on Clement explicitly in relation to the Gnostics, which dimension of Clement had been previously ignored (perhaps because there are so many others?).

The body of Clement's work is saturated (also) with Jewish\footnote{\url{http://www.bgbucur.com/PDFuri/ClementHierarchy.pdf}} mythology.

It is unnecessary to rehearse his spiritual lineage through Papias to the apostles, or to trace his influence in later Church mystical theology, but Dionysius was an indirect pupil, and Dionysius' writings are arguably ``unorthodox'' enough to accommodate a possible Hermetic root source. The notion of the deep ``I'' emanating and proceeding and then returning\footnote{\url{http://www.reversespins.com/dionysius.html}} to God are of particular import. His Stoical influences are so far beyond doubt as to be unremarkable.

We can conclude that Clement's ``canon'' included all things because ``all things were from God''. Thus, paganism cannot evaporate or be subordinated:

\begin{quotex}
In fact, nevertheless the thief possesses really, what he has possessed himself of dishonestly, whether it be gold, or silver, or speech, or dogma. The ideas, then, which they have stolen, and which are partially true, they know by conjecture and necessary logical deduction: on becoming disciples, therefore, they will know them with intelligent apprehension.
\end{quotex}


Here, Clement virtually argues that Christian Gnosis in fact re-illuminates that which it has already illuminated indirectly, through union with its derivatives. This return and reconciliation (similar to the doctrine of recapitulation given in Colossians) of paganism and Christianity should make clear that it was possible to achieve a balance between the two which allowed ``each to be for the other, but differently''.

Clement himself admitted that there was a ``secret Gospel'', but claimed that it was not what the Gnostics pretended it to be. Might it have looked more like Mouravieff's\footnote{\url{http://www.bibliotecapleyades.net/archivos_pdf/mouravieff3.pdf}} work?

Even his pupil Origen's echo is true to this presupposition \& Clement's teaching:

\begin{quotex}
Jesus conversed with His disciples in private, and especially in their sacred retreats, concerning the Gospel of God; but the words which He uttered have not been preserved, because it appeared to the evangelists that they could not be adequately conveyed to the multitude in writing or in speech... and they saw... what things were to be committed to writing, and how this was to be done, and what was by no means to be written to the multitude, and what was to be expressed in words, and what was not to be so conveyed. \flright{\textit{Contra Celsus}, VI, 18}
\end{quotex}


So that Clement's citation of Paul's words make perfect sense:

\begin{quotex}
Take also the Hellenic books, read the Sibyl, how it is shown that God is one, and how the future is indicated. \flright{\textsc{Clement}, \textit{Stromata}, Book VI, chap. VII}
\end{quotex}


God is one, truth is one. It should be clear by now that Titus Flavius Clemen's thought-world (still very similar to St. Paul's) was a far cry from that of Tertullian, Cyprian, Iranaeus, or even Augustine, and that the Latin world's loss of the Alexandrian library was a very symbolic event. Clement seems to have been a point of intersection for an amazing amount of eclectic influence, and it would be very surprising if that were merely a matter of chance.


\flrightit{Posted on 2011-04-15 by Logres}

\begin{center}* * *\end{center}

\begin{footnotesize}
\begin{sffamily}

\texttt{Logres on 2011-04-16 at 11:16 said:}

Clement is not arguing at all that Christianity ``saved'' paganism -- that is an Augustinian argument.

\hfill

\texttt{Spock on 2011-04-17 at 10:03 said:}

Whether completed or saved is splitting hairs. My reference to saving was to the gentiles not to paganism where I used Clement's word ``completing.'' However, salvation and completion you must admit carry an almost identical value.

\hfill

\texttt{Logres on 2011-04-17 at 19:35 said:}

Spock, they could mean the same thing, but Clement's ``mental and spiritual furniture'' indicate they don't. Clement does not think that the Gospel ``saves'' the pagan schools. And it completes them only by essentially being the same thing, \textbf{except more so}.

It might be worth while pondering the effect Augustine and his milieu had on the infant Church. Although I might be willing to concede that Saint Paul comes down on the side of ``Platonism for the masses'' (not my original phrase), at least in general effect, it seems clear that the collapse of the Empire, the persecutions (which annihilated true apostolic succession in a lineal sense), and the general confusion of the Late Roman period combined to form in Saint Augustine a decisive favoring of Latinate ``City of God'' style religion as opposed to authentic and alternate streams present in the Church, which were still there as late as Boethius. Even the Orthodox want to credit Augustine here, for his insistence (echoed today by V. Lossky, the premier mystical theologian)that there is no such thing as ``hidden'' wisdom apart from dogma. Clement's position is far more subtle. While he agrees (against the Gnostics) that dogma is inseparable from esoteric truth, he would agree (against Augustine and others) also that esoteric truth must be maintained to keep dogma intact. This retention of esoteric truth, furthermore, for Clement is not passed down through open writing, or even open proclamation, but the presence within the Church of ``those who know''. ``Those who know'' are educated, or self-educated, by the wisdom available to all those who are receptive to it, an inner truth which is embodied in diverse traditions and schools and wisdoms known today as ``paganism''. Clement's inferred position from this is that a Christianity divorced from such is either no true Christianity, or a Christianity unwilling to accept education in the school of wisdom, and therefore at risk for being superseded in the same manner that the Jews (presumably) were superseded when they arrogated perfection to their exoteric assemblies.

\hfill

\texttt{Logres on 2011-04-17 at 19:37 said:}

For example,\url{http://members.ozemail.com.au/~moorea/agapetae.html}

\begin{quotationx}\footnotesize

It seems that there was in the first full rush of the Church, an attempt, encouraged by the Apostles, to `sublimate'. But the experimenters did not call it that. The energy of the effort was in and towards the crucified and Glorified Redeemer, towards a work of exchange and substitution, a union on earth and in heaven with the love which was now understood to be capable of loving and being loved. In some cases it failed. But we know nothing -- most unfortunately -- of the cases in which it did not fail, and that there were such cases seems clear from St Pauls quite simple acceptance of the idea. By the time of Cyprian, Bishop of Carthage in the 3rd Century, the ecclesiastical authorities were more doubtful. The women, sub-introductae as they were called -- apparently slept with their companions without intercourse; Cyprian does not exactly disbelieve them, but he discourages the practice. And the synod of Elvira (305) and the Council of Nicea (325) forbade it altogether. The great experiment had to be abandoned because of Scandal.

Tolstoy put the crude objection in the \textit{Kreutzer Sonata}, and Cyprian more or less agreed. ``But then excuse me, why do they go to bed together?'' Both wise men were justified as against a great deal of sentimental lust and sensual hypocricy. But even Cyprian and Tolstoy did not understand all the methods of the blessed spirit in Christendom. The prohibition was natural. Yet it seems a pity that the Church, which realised once that she was founded on a Scandal, not only to the world, but to the soul, should be so nervously alive to scandals. It was one of the earliest triumphs of ``the weaker brethren,'' those innocent sheep who by mere volume of imbecility have trampled over many delicate flowers in Christendom. ...

\end{quotationx}

\end{sffamily}
\end{footnotesize}

